\section{Empirical Setup}
\label{sec:setup}

% We evaluate three research questions:

% \begin{itemize}
%     \item \textbf{RQ1: What capabilities does model merging compose?}
%     We test whether merging combines transform-specific adaptations or more
%     general program-reasoning capabilities.

%     \item \textbf{RQ2: How effectively can capabilities learned from individual
%     obfuscations be composed for stacked obfuscation?}
%     We compare prompting, pooled training, consistency training,
%     inference-time routing, and parameter-space merging.

%     \item \textbf{RQ3: How robust and generalizable is the capability produced
%     by model merging?}
%     We evaluate deeper stacks, an unseen obfuscation family, cross-language
%     transfer, and reversed input/output reasoning.
% \end{itemize}


\subsection{Models}

We evaluate eight instruction-tuned models spanning four model families:
CodeLlama-7B, CodeLlama-13B, and CodeLlama-34B~\cite{code_llama};
Llama-3.1-8B~\cite{llama3herd};
StarCoder2-15B~\cite{lozhkov2024starcoder2stackv2};
Gemma-3-12B~\cite{gemma3};
CodeGemma-7B~\cite{codegemma}; and
Granite-3.1-8B~\cite{granite31}.
The panel includes both general-purpose and code-specialized models from
7B to 34B parameters. We additionally use the pretrained Llama-3.1-8B
checkpoint for the one-shot baseline. All adapted models use rank-32 LoRA
with the same target modules.


\subsection{Dataset}

Our Python corpus contains 2,231 programs drawn from
APPS~\cite{hendrycks2021apps},
CRUXEval~\cite{cruxeval}, and
HumanEval~\cite{Chen2021}.
We retain 1,584 APPS programs, 543 CRUXEval programs, and all 104
HumanEval programs. We split the corpus by program so that no transformed
variant of a training program appears in evaluation. Each program is evaluated
on three fixed inputs, with reference outputs obtained by executing the clean
program.

\begin{table}[t]
\centering
\small
\caption{\textbf{Dataset summary.} Corpus composition and program-level split.}
\label{tab:dataset}
\setlength{\tabcolsep}{5pt}
\renewcommand{\arraystretch}{1.05}

\begin{tabular}{@{}lr@{\hspace{18pt}}lr@{}}
\toprule
\multicolumn{2}{c}{\textbf{Corpus}} &
\multicolumn{2}{c}{\textbf{Split and properties}} \\
\cmidrule(r){1-2}
\cmidrule(l){3-4}

\textbf{Dataset} & \textbf{Programs} &
\textbf{Property} & \textbf{Value} \\
\midrule

APPS~\cite{hendrycks2021apps} & 1,584 &
Training programs & 1,563 \\

CRUXEval~\cite{cruxeval} & 543 &
Validation programs & 111 \\

HumanEval~\cite{Chen2021} & 104 &
Held-out programs & 557 \\

\textbf{Total} & \textbf{2,231} &
Inputs per program & 3 \\

& &
Held-out clean instances & 1,671 \\

& &
Mean program length & 8.4 LOC \\

& &
Maximum program length & 57 LOC \\

\midrule

JavaScript & 168 programs &
JavaScript instances & 504 \\

\bottomrule
\end{tabular}
\end{table}

Coverage varies by transformation because some structural and encoding
transforms are not applicable to every program. All paired comparisons are
therefore computed over the intersection of programs available in the
conditions being compared.


\subsection{Obfuscation Types}

We organize the obfuscation conditions into \textbf{seen single transforms},
\textbf{held-out encoding transforms}, and \textbf{stacked composites}.
See Table~\ref{tab:taxonomy} for the full definitions and coverage of each
condition.

\paragraph{Seen single transforms.}
The main training set contains five obfuscations:
three identifier transforms---\texttt{L1b}, \texttt{L1r}, and \texttt{L2}---and
two structural transforms---\texttt{S1} and \texttt{S2}.
\texttt{L1b} performs adversarial identifier renaming, \texttt{L1r} performs
random renaming, and \texttt{L2} minifies identifiers.
\texttt{S1} applies control-flow flattening, while \texttt{S2} combines opaque
predicates with dead-code insertion.
Each specialist is trained on exactly one of these conditions, and no specialist
is trained on a stacked obfuscation. Together with the clean-code
condition \texttt{L0}, these form the six specialist conditions used by the
default breadth and merge systems.

\paragraph{Structural components.}
We additionally separate \texttt{S2} into its two constituent mechanisms:
\texttt{S3} contains dead-code insertion only, and \texttt{S4} contains opaque
predicates only. These conditions are used to construct and analyze stacked
composites; they are not part of the five-transform training set used by the
breadth and merge systems.

\paragraph{Held-out encoding family.}
The encoding family is excluded from the main systems' training data.
\texttt{X1} combines string encoding with semantics-preserving arithmetic
rewriting. \texttt{X1m} contains only the arithmetic component, and
\texttt{X1s} contains only the string-encoding component.
\texttt{H1} is a quarantined held-out obfuscator using a distinct
Base64/RC4-based string encoding together with mixed-boolean-arithmetic
rewriting. Thus, \texttt{X1}, \texttt{X1m}, and \texttt{X1s} define the
held-out encoding conditions used in the main generalization experiments,
while \texttt{H1} is reserved as the quarantined reference family.

\paragraph{Stacked composites.}
Stacked conditions apply multiple transforms sequentially to the same clean
program, with order treated as part of the condition because later transforms
may modify code introduced by earlier ones. We evaluate sixteen stacks:
six depth-2 stacks composed only of seen transforms
(\texttt{L1b}, \texttt{L1r}, \texttt{L2}, \texttt{S1}, \texttt{S2},
\texttt{S3}, or \texttt{S4});
four deeper stacks of seen transforms at depth 3 or 4; and six stacks that
contain the held-out encoding family through \texttt{X1}, \texttt{X1m}, or
\texttt{X1s}. No system is trained directly on any stacked condition.


\begin{table*}[t]
\centering
\scriptsize
\setlength{\tabcolsep}{3.8pt}
\renewcommand{\arraystretch}{1.10}
\caption{\textbf{Obfuscation conditions.}
Single transforms define the atomic conditions used for training and
generalization evaluation; stacked conditions are constructed only for
evaluation.}
\label{tab:taxonomy}

\begin{tabular}{@{}llp{0.42\textwidth}cl@{}}
\toprule
\textbf{Code / group} &
\textbf{Family} &
\textbf{Transformation} &
\textbf{Programs} &
\textbf{Role} \\
\midrule

\multicolumn{5}{@{}l}{\textbf{Clean reference}} \\
L0 &
none &
Normalized source with comments and docstrings removed &
557 &
clean \\

\midrule
\multicolumn{5}{@{}l}{\textbf{Seen single-transform conditions}} \\

L1b &
identifier &
Adversarial renaming of bindings and the entry function &
553 &
train / eval \\

L1r &
identifier &
Random hexadecimal renaming of all bindings &
557 &
train / eval \\

L2 &
identifier &
Sequential identifier minification and removal of type annotations &
557 &
train / eval \\

S1 &
structural &
Control-flow flattening into a dispatch loop with randomized states &
416 &
train / eval \\

S2 &
structural &
Opaque predicates combined with never-called dead helper functions &
556 &
train / eval \\

\midrule
\multicolumn{5}{@{}l}{\textbf{Structural components used in stacks}} \\

S3 &
structural &
Dead-code insertion using never-called helper functions &
557 &
stacks \\

S4 &
structural &
Opaque predicates inserted at statement boundaries &
556 &
stacks \\

\midrule
\multicolumn{5}{@{}l}{\textbf{Held-out encoding family}} \\

X1 &
encoding &
String encoding and semantics-preserving arithmetic rewriting &
405 &
held-out eval \\

X1m &
encoding &
Arithmetic-rewriting component of X1 &
351 &
held-out eval \\

X1s &
encoding &
String-encoding component of X1 &
246 &
held-out eval \\

H1 &
encoding &
Base64/RC4 string encoding with mixed-boolean-arithmetic rewriting &
-- &
quarantined \\

\midrule
\multicolumn{5}{@{}l}{\textbf{Stacked evaluation conditions}} \\

Seen depth-2 &
composite &
Two sequentially applied seen transforms &
varies &
6 conditions \\

Seen depth-3/4 &
composite &
Three or four sequentially applied seen transforms &
varies &
4 conditions \\

Held-out stacks &
composite &
Stacks containing X1, X1m, or X1s together with seen transforms &
varies &
6 conditions \\

\bottomrule
\end{tabular}
\end{table*}


\subsection{Systems}

All trainable systems use rank-32 LoRA adapters with $\alpha=64$ and identical
target modules. Systems differ only in their training data or composition
mechanism as described in Sections \ref{sec:merging} and \ref{sec:baselines}.


\subsection{Measurement}

Let $p_i$ denote a program, $x_i$ an input, and
$y_i = p_i(x_i)$ its reference output.

\paragraph{Seed noise floor.} A bootstrap over programs does not contain training variation, so
we measure that separately. Three independently seeded runs of the same recipe --- seeds 17, 101
and 202, six specialists and pooled training, on CodeLlama-7B --- give a median range of $0.66$
accuracy points across 21 arm $\times$ condition-group combinations, with a maximum of $2.39$
(Table~\ref{tab:seednoise}). The noise is not uniform: it is roughly three times wider on the
held-out family (median $1.57$) than on clean code (median $0.66$). The two quantities answer
different questions and neither subsumes the other, so we do not narrate single-seed differences
below roughly one point on clean code or the seen singles, or below roughly $2.4$ points on the
held-out family. Measured on one model; the other seven are not measured and may differ.

% GENERATED by scripts/analysis/83_seed_variance.py -- do not edit by hand.
\begin{table}[t]
\centering\small
\caption{\textbf{The seed noise floor.} Range (max $-$ min) in accuracy points across three independently seeded runs of the same recipe --- seeds 17, 101 and 202 --- on CodeLlama-7B. \textbf{The noise is not uniform}: it is roughly three times wider on the held-out family than on clean code, and that is the condition the paper's sharpest comparisons are made on. These are \emph{training} variation and are disjoint from the program-clustered bootstrap intervals reported elsewhere, which measure sampling variation over programs and contain no seed effect at all. A narrated difference smaller than the relevant entry is not distinguishable from a re-rolled seed. Measured on one model; other models are not measured and may differ.}
\label{tab:seednoise}
\begin{tabular}{@{}lrrr@{}}
\toprule
arm & \texttt{L0} & singles & unseen (\texttt{X1}) \\
\midrule
\texttt{L0} & 0.60 & 0.66 & 0.58 \\
\texttt{L1b} & 1.44 & 0.28 & 2.39 \\
\texttt{L1r} & 0.24 & 0.90 & 1.57 \\
\texttt{L2} & 0.66 & 0.53 & 2.31 \\
\texttt{S1} & 0.06 & 0.63 & 0.74 \\
\texttt{S2} & 1.32 & 0.30 & 1.65 \\
\emph{breadth} & 1.50 & 1.44 & 0.08 \\
\midrule
\textbf{median} & \textbf{0.66} & \textbf{0.63} & \textbf{1.57} \\
\bottomrule
\end{tabular}
\end{table}


\paragraph{Forward accuracy.}
Given $(p_i, x_i)$, the model predicts an output $\hat{y}_i$. We measure
normalized exact-match accuracy:
\[
\mathrm{Acc}_{\mathrm{fwd}}
=
\frac{1}{N}
\sum_{i=1}^{N}
\mathbf{1}
\!\left[
\operatorname{norm}(\hat{y}_i)
=
\operatorname{norm}(y_i)
\right].
\]

\paragraph{Backward accuracy.}
Given $(p_i, y_i)$, the model predicts an input $\hat{x}_i$. Because multiple
inputs may produce the same output, correctness is determined by execution:
\[
\mathrm{Acc}_{\mathrm{bwd}}
=
\frac{1}{N}
\sum_{i=1}^{N}
\mathbf{1}
\!\left[
p_i(\hat{x}_i) = y_i
\right].
\]
We additionally report \textbf{reference-input recovery}, the fraction of
predictions that exactly recover the reference input $x_i$.

\paragraph{Statistical comparisons.}
All system comparisons are paired over the intersection of programs available
in the compared conditions. We report confidence intervals from a
program-clustered bootstrap with 2,000 resamples. 